% !TeX root = ../../Book3.tex
\section{单值性}
\label{b3:sec:0903}

单值性定理不声称每个芽都能延拓。它要求沿指定区域内所有路径可延拓，然后以同伦控制路径依赖。在单连通区域中，这种依赖完全消失。

\subsection{同伦不变性}
\label{b3:subsec:090301}

\begin{theorem}\label{b3:thm:09-homotopy-continuation}
设初始芽可沿区域 \(U\) 内从 \(a\) 出发的每条路径延拓。端点固定同伦的两条路径给相同末芽。
\end{theorem}
\begin{proof}
设 \(H(s,t)\) 为同伦。对每个 \(s_0\)，路径 \(H(s_0,\cdot)\) 有延拓；有限链稳定性和 \(H\) 的一致连续性给 \(s_0\) 的一个邻域，其中所有路径使用相同有限代表，末芽相同。因此末芽随 \(s\) 局部恒定。参数区间连通，故从0到1始终相同。
\end{proof}
条件“每条路径可延拓”在证明中用于每一个 \(s_0\)。若同伦扫过一个分支奇点，某个中间路径的延拓不存在，局部恒定论证就断开；不能只核对最初与最后两条路径都能延拓。

\ProseParagraphBreak
有限链稳定性需要端点固定，否则“末芽相同”没有共同中心可以比较。若只把端点轻微移动，正确结论是末芽由同一个终端代表在各自中心给出。端点固定同伦恰使这些终端代表落在同一个芽集合里，局部恒定才有意义。自由同伦的闭路若基点移动，monodromy 作用要先沿基点移动路径作共轭识别；这与基本群换基点的规则一致。

\ProseParagraphBreak
对数给一个可以手工核验的同伦过程。路径沿小盘移动时，每个局部对数的差为 \(2\pi\mathrm{i}\) 整数，连续变形不能突然改变该整数，所以绕数固定。只要同伦避开0，末分支始终一致；若允许同伦穿过0，一个中间路径没有对数芽，整数便可以在这一失效处改变。数学陈述中应把“在 \(U\) 中同伦”写全，不能只说两个图形可以在整个平面拉成一样。

\subsection{单值性定理}
\label{b3:subsec:090302}

\term[dan-zhi-xing-dingli]{单值性定理}[monodromy theorem]由同伦不变性消除单连通域上的路径依赖。

\begin{theorem}[单值性定理]\label{b3:thm:09-monodromy}
设 \(U\) 单连通，初始芽可沿 \(U\) 内每条从中心出发的路径延拓。则所有延拓芽由 \(U\) 上同一个全纯函数给出。
\end{theorem}
\begin{proof}
给 \(z\in U\)，选一条从 \(a\) 到 \(z\) 的路径，并以末芽的值定义 \(F(z)\)。单连通性使任意两条这样的路径端点固定同伦，前一定理保证定义不依赖选择。固定 \(z\) 后，末代表在小盘内全纯；把到 \(z\) 的路径再接一条小盘内线段，就给附近点的值，故 \(F\) 在该盘等于这个代表。因此 \(F\) 全纯且具有所有所需芽。
\end{proof}
若区域不单连通，延拓仍由路径同伦类控制，闭路类作用在可达芽上，称为单值群的作用或 monodromy 作用。它可以平凡，例如单值有理函数，也可以非平凡；基本群非平凡不强迫每个函数都多值。

\begin{proposition}[线性微分方程的延拓]\label{b3:prop:09-linear-ode}
设 \(a,b\) 在区域 \(U\) 全纯，\(z_0\in U\)、\(c\in\mathbb C\)。初值问题
\[
 y'+a(z)y=b(z),\qquad y(z_0)=c
\]
在 \(z_0\) 有唯一全纯解芽，并可沿 \(U\) 内每条路径延拓。若 \(U\) 单连通，这些芽由同一个全纯函数给出。
\end{proposition}
\begin{proof}
小盘内选 \(A'=a\)、\(A(z_0)=0\)，再选 \(B'=e^A b\)、\(B(z_0)=0\)。则
\[
 y=e^{-A}(c+B)
\]
直接满足方程与初值。两个解之差乘 \(e^A\) 后导数为0，因此唯一。给定路径的紧像可用有限圆盘链覆盖；在每段上取局部积分因子，以前一段重叠处的值为新初值继续。局部表达式在整个圆盘全纯，有限链不会出现有限步的解爆破。唯一性使重叠代表相等，故得到沿全路径的延拓。单连通时应用单值性定理即得全局函数。
\end{proof}
这里线性性确实保证每个圆盘上任意初值都可解到整个盘。非线性方程未必如此。例如 \(y'=y^2\)、\(y(0)=1\) 的解为 \(1/(1-z)\)，即使方程系数在整个平面全纯，解也在1出现极点。因此“系数全纯”本身不够保证所有解沿全部路径以全纯芽延拓。

\ProseParagraphBreak
在非单连通域，线性方程的闭路作用是仿射变换。沿路径续接积分因子后，\(A\) 在闭路返回时增加周期 \(\alpha=\int a\,\mathrm{d}z\)。于是齐次解 \(e^{-A}\) 乘 \(e^{-\alpha}\)；非齐次项另给一个与初值无关的常数，末初值形如 \(c\mapsto e^{-\alpha}c+\beta\)。两条闭路的作用按仿射变换复合，它们可能不交换。对于 \(y'=1/z\)，作用是平移 \(c\mapsto c+2\pi\mathrm{i}n\)；对于 \(y'+(\nu/z)y=0\)，作用是乘 \(e^{-2\pi\mathrm{i}\nu n}\)。前者没有单值解，后者在 \(\nu\in\mathbb Z\) 时所有解单值，在 \(\nu\notin\mathbb Z\) 时零解仍单值。单值性可以取决于初值，而不只是方程。

\subsection{多值函数}
\label{b3:subsec:090303}

“多值函数”可以精确地理解为一张芽曲面上的单值全纯函数，同时投影到底层区域。沿闭路延拓改变的是位于同一底层点上的芽，不是一个函数在同一空间点取多个值。对 \(z^{1/m}\)，绕0一周乘 \(e^{2\pi\mathrm{i}/m}\)，故有 \(m\) 个芽；对 \(\log z\)，绕行产生无限多个相差 \(2\pi\mathrm{i}n\) 的芽。

\ProseParagraphBreak
有些函数确有自然边界。例如 \(F(z)=\sum_{n\ge0}z^{2^n}\) 在单位盘全纯，单位圆上每个 \(2^m\) 次单位根 \(\zeta\) 都是不可延拓点。当 \(r\uparrow1\)，从 \(n\ge m\) 起有 \((r\zeta)^{2^n}=r^{2^n}>0\)，其和趋于无穷：给定任意有限项数，先令这些项都接近1，再令项数增加。前面有限项保持有界，所以 \(F(r\zeta)\) 无界。这些根在圆周稠密；若能跨某一圆周点延拓到小盘，该盘与圆周的一段开弧内就不应有这样的奇点，矛盾。单值性定理不能跨越没有延拓的自然边界。

\ProseParagraphBreak
把芽的全部闭路作用汇集起来，可以看到函数与基本群之间的信息损失。某些不同闭路作用相同：平方根只记录绕数的奇偶，对数只记录一个整数，而双孔域的基本群还保留非交换次序。一个群元素对初始芽作用平凡，也未必对全部可达芽作用平凡；要定义整个芽纤维上的置换作用，必须追踪每个芽的继续延拓。后面的覆盖 monodromy 则以完整纤维为对象，每条闭路给该纤维的一个双射。

\ProseParagraphBreak
自然边界例子的结论也强于“Taylor 级数不收敛到外面”。在任意一个圆周点延拓，代表的小盘会碰到一个开圆弧；稠密的不可延拓点落在该弧内，限制回原盘会与原级数相等，因而也会延拓那个已知奇点，矛盾。这里恒等定理负责把新代表与原函数接起来，稠密性负责把局部延拓变成对某个具体奇点的矛盾，两项缺一不可。

\subsection{\texorpdfstring{\(\log z\)}{log z}、\texorpdfstring{\(z^\alpha\)}{z 的 \AnalysisPdfAlpha 次幂} 与反函数}
\label{b3:subsec:090304}

对 \(\log z\)，末值变化为 \(2\pi\mathrm{i}\operatorname{Ind}(\gamma,0)\)。对 \(z^\alpha=\exp(\alpha\log z)\)，绕行 \(n\) 次乘 \(e^{2\pi\mathrm{i}\alpha n}\)。若 \(\alpha=p/q\) 为最简有理数，有 \(q\) 层；若 \(\alpha\) 非有理实数或非实复数，层数无限。非实 \(\alpha\) 的因子模也改变，不能将所有分支都画成同一个模长的旋转。

\ProseParagraphBreak
反函数芽在导数非零处由局部逆函数定理存在。指数的逆芽就是对数；多项式的逆在临界值之外可沿路径延拓，闭路可能置换根。例如 \(z^2\) 的逆绕0交换两根。对于 \(z^3-3z\)，临界点为 \(\pm1\)，临界值为 \(\mp2\)；绕这些值分别交换相遇的两个根，两个交换一般不交换。根的全部对称函数仍单值，说明多值根与单值系数可以同时存在。

\ProseParagraphBreak
寻找反函数延拓的障碍时，临界值是局部障碍，但对非多项式整函数还可能出现有限渐近值。\(e^z\) 没有临界点，却不能在0附近取得逆芽；沿实轴向负无穷的像趋于0。不能仅从导数处处非零就宣布它在整个目标平面是一张覆盖。

\ProseParagraphBreak
双平方根的例子进一步区分单个分支点与整体绕行。对 \(\sqrt{z^2-1}\)，在 \(\mathbb C\setminus\{-1,1\}\) 中绕1或绕 \(-1\) 一次都使根变号；绕两点的闭路使 \(z^2-1\) 绕0两次，因而根不变。可以在连接 \(-1\) 与1的线段外选一个全纯分支，并以 \(F(z)/z\to1\) 在无穷远附近确定它。另一个分支为 \(-F\)。两根的平方和系数始终单值，不代表每个根都单值。

\ProseParagraphBreak
还要区分函数值回来和芽回来。一个闭路延拓后，末值偶然等于起值，末导数却可能不同；判定单值应比较末芽。对于多项式逆芽，起点取一个正则值，不同根本身不同，所以值足以区分这些逆芽；对于一般延拓问题没有这种保证。若在一个有聚点的小邻域中全部值回来，恒等定理才把值的比较提升为芽的比较。

\ProseParagraphBreak
整函数的有限渐近值还给逆芽的非紧致障碍。\(\int_0^z e^{-t^2}\,\mathrm{d}t\) 的导数处处非零，但沿正实轴趋于有限极限。某个逆分支沿像曲线向这个极限逼近时，源点逃向无穷，不能在该极限处接成有限源点的局部逆。局部逆函数定理没有失效；失效的是把局部逆图串联成整个目标上的完整覆盖所需的全局控制。多项式没有有限源逃逸，因为其适当性由最高次项增长保证，下一章覆盖论将利用这点。

\ProseParagraphBreak
割线选择可以改变一个分支的定义域，却不改变全部延拓的芽曲面。\(\log z\) 可在负实轴外定义主分支，也可沿任一从0通向无穷远的合适简单割线定义一个分支；这些单值域各自只选择同一无限层曲面上的一块。割线两侧的值相差 \(2\pi\mathrm{i}\)，是所选局部截面在重叠处的差，并非函数在曲面上不连续。画图时割线是为选分支而引入，未必是原函数本身的一条自然奇性边界。

\ProseParagraphBreak
对于 \(z^\alpha\)，若 \(\alpha=p/q\) 约分后有 \(q\) 层，其非零芽曲面可用参数 \(t\) 表示 \(z=t^q,F=t^p\)。若 \(p<0\)，\(t=0\) 处 \(F\) 为极点；若 \(p>0\)，为零点。底层映射在那里分歧，取决于 \(q\)，而函数本身的零极点阶取决于 \(p\)。分支重数与函数值阶数不是同一个整数，这在代数曲线的局部坐标中尤其重要。
